% !TeX root = ../navier-stokes-manuscript.tex
% ============================================================
% 2.6 Finite Adjacent Rectangles
% ============================================================

\subsection{Finite Adjacent Rectangles}\label{sub:FiniteAdjacentRectangles}

Fix a temporal depth \(N\in\Nzero\), a spatial height parameter \(M\in\Nzero\), and a strict interval \(I=(0,T)\) with \(T<T_*\).
The temporal rows will be
\[
  V_n:=\partial_t^n u,
  \qquad
  0\leq n\leq N+1.
\]
The correct spatial pair around the middle height \(H^M\) is
\[
  H^{M+1}(\R^3)
  \hookrightarrow
  H^M(\R^3)
  \hookrightarrow
  H^{M-1}(\R^3).
\]
For \(M=0\), the lower space is \(H^{-1}\) in the standard inhomogeneous Sobolev scale.

This adjacent pair has a specific temporal job.
If
\[
  V_n\in L^2(I;H^{M+1}),
  \qquad
  \partial_tV_n=V_{n+1}\in L^2(I;H^{M-1}),
\]
then the standard Hilbert-triple trace mechanism gives
\[
  V_n\in C(\overline I;H^M).
\]
The upper row controls one spatial derivative above the trace height, and the time derivative is allowed to live one spatial derivative below it.
Their two-order separation matches the parabolic scaling already visible in \(\partial_t-\nu\Delta\).
The middle \(H^M\) value is the trace paid for by two adjacent spatial rows.

The first useful example takes \(N=1\) and \(M=2\).
We read \(u\) and \(\partial_tu\) at the upper spatial height \(H^3\), then read their time derivatives \(\partial_tu\) and \(\partial_t^2u\) at the lower height \(H^1\).
The resulting finite quantity is
\begin{align*}
  \mathfrak R_{1,2}(u;T)
  :={}&
  \norm{u}_{L^2(0,T;H^3)}^2
  +
  \norm{\partial_tu}_{L^2(0,T;H^3)}^2
  \\
  &+
  \norm{\partial_tu}_{L^2(0,T;H^1)}^2
  +
  \norm{\partial_t^2u}_{L^2(0,T;H^1)}^2.
\end{align*}
The first and third terms form the adjacent pair for \(u\), so they give \(u\in C([0,T];H^2)\).
The second and fourth terms form the adjacent pair for \(\partial_tu\), so they give \(\partial_tu\in C([0,T];H^2)\).
The repeated \(\partial_tu\) is the join between these two blocks.
Its upper \(H^3\) membership already implies its lower \(H^1\) membership, while recording both appearances keeps its two structural roles visible.

The same example can be read as a small rectangle of the mixed jet:
\[
  \begin{array}{c|cc}
  & n=0 & n=1 \\
  \hline
  \text{upper row }H^3
  & u & \partial_tu \\
  \text{lower derivative row }H^1
  & \partial_tu & \partial_t^2u
  \end{array}
\]
The horizontal coordinate is temporal depth, and the vertical coordinate is the adjacent spatial pair.
The word rectangle names this finite coordinate pattern.
It is not a physical box inside the fluid.

Nothing in the construction depends on the particular cutoffs \(1\) and \(2\).
At arbitrary finite depths, define
\begin{align*}
  \mathfrak R_{N,M}(u;T)
  :={}&
  \sum_{n=0}^{N}
  \norm{\partial_t^n u}_{L^2(0,T;H^{M+1})}^{2}
  \\
  &+
  \sum_{n=0}^{N}
  \norm{\partial_t^{n+1}u}_{L^2(0,T;H^{M-1})}^{2}.
\end{align*}
Each index \(n\) contributes one field row at upper spatial height and its time derivative at lower spatial height.
Consequently, every \(0\leq n\leq N\) has the middle trace
\[
  \partial_t^n u
  \in
  C\bigl([0,T];H^M(\R^3)\bigr).
\]
Increasing \(N\) exposes more temporal responses, while increasing \(M\) raises the spatial height of every exposed row.
The two cutoffs remain independent.

The norm measures the velocity coordinates of a finite view.
The pressure-complete record must also retain the corresponding pressure rows and the differentiated equations that connect them.
Write \(\cR_{N,M}(u,p;T)\) for the record consisting of the velocity rows counted by \(\mathfrak R_{N,M}(u;T)\), the pressure coordinates determined by
\[
  -\Delta\Pi_{m,\alpha}
  =
  \partial_i\partial_j
  \partial_t^m\partial_x^\alpha(u_i u_j),
\]
and every mixed recurrence needed inside the chosen finite depths.
Attaching these equations prevents a list of Sobolev functions from being mistaken for a \NavierStokes\ derivative record.
The finite view still carries viscosity, pressure, incompressibility, transport, and the binomial interactions among its coordinates.

One rectangle pays only a finite part of smoothness.
For fixed \(N\) and \(M\), it gives the middle traces available at those cutoffs.
The all-orders intersection arises by allowing every finite \((N,M)\), with a separate finite bound permitted at each coordinate.
Those separate views can form one intersection only when their shared rows agree.
